\section{Consequences and limitations}\label{sec:consequences}
For a nonzero real solution, the weight in \eqref{eq:4-4} is strictly positive, so v and f change sign. Since k vanishes near 0, \(v=-\mathsf Tv\) there. If f vanished on a terminal-time neighborhood, the entire function \(\mathsf Tv\) would vanish identically. Equation \eqref{eq:2-3} and density of polynomials in H would give \(v=0\), a contradiction. The same relation makes f analytic near both endpoints, but f cannot be real analytic on a connected neighborhood of \([0,T]\): otherwise k would be analytic there and vanish identically because it is zero near 0.

The entire function \(G(z)=\int_0^T t e^{zt}v(t)dt\) is nonzero for nonzero v by polynomial density. Its zeros cannot accumulate at 0; since \(\delta_n\to0\), \eqref{eq:4-5} shows that all but finitely many memory coefficients are nonzero. No prescribed Gevrey order, numerical closed form, quantitative control cost, nonlinear preparation theorem, or full-system non-identifiability follows. The theorem's proof and its source applications have been audited only by the same agent; independent mathematical and physical-model review, and chronological priority assessment, remain necessary.
